\section{Methodik} 

In diesem Abschnitt widmen wir uns den angewandten Methoden unserer Arbeit und erläutern die verwendeten Modelle sowie ihre Metriken und Algorithmen.

Die Visualisierung \ref{fig:07-simplified_overview} stellt unsere Arbeit einfach dar und gibt einen Überblick, um ein grundlegendes Verständnis dafür zu vermitteln, worum es in der Arbeit geht.

\begin{figure}[H]
    \centering
    \includegraphics[width=13.5cm]{01-images/04-methodik/simplified_overview.png}
    \caption{Simplifizierte Version unseres Vorgehens}
    \label{fig:07-simplified_overview}
\end{figure}

Eine weiterführende, detailreichere Visualisierung ist \ref{fig:06-overview}, die die oben gezeigte Visualisierung weiter aufschlüsselt und genauer beschreibt.

\begin{figure}[H]
    \centering
    \includegraphics[width=\linewidth]{01-images/04-methodik/overview.png}
    \caption{Methodik Übersicht}
    \label{fig:06-overview}
\end{figure}


\subsection{Technische Umsetzung}

Für die technische Umsetzung unserer Bachelorarbeit haben wir auf bewährte und bekannten Data Science Frameworks zurückgegriffen. Darunter gehören: Numpy, Pandas, Matplotlib, \& Seaborn. Für die Modelle sowie die Verarbeitung von Bildern nutzen wir PyTorch sowie die Erweiterung PyTorch-Lightning und tracken den Trainingsprozess sowie unsere Modelle in Weights \& Bias.
